\documentclass[conference]{IEEEtran}
\IEEEoverridecommandlockouts
\usepackage{cite}
\usepackage{amsmath,amssymb,amsfonts}
\usepackage{algorithmic}
\usepackage{graphicx}
\usepackage{textcomp}
\usepackage{xcolor}
\def\BibTeX{{\rm B\kern-.05em{\sc i\kern-.025em b}\kern-.08em
    T\kern-.1667em\lower.7ex\hbox{E}\kern-.125emX}}

\begin{document}

\title{Typing Tutor: A High-Performance OS-Inspired System with Modular C Engine and Full-Stack Integration}

\author{\IEEEauthorblockN{Siddhartha Shukla}
\IEEEauthorblockA{\textit{Enrollment No: 230104} \\
\textit{Rishihood University}}
\and
\IEEEauthorblockN{Sugat Athawale}
\IEEEauthorblockA{\textit{Enrollment No: 230109} \\
\textit{Rishihood University}}
\and
\IEEEauthorblockN{Piyush Kaushal}
\IEEEauthorblockA{\textit{Enrollment No: 230112} \\
\textit{Rishihood University}}
}

\maketitle

\begin{abstract}
This paper presents the design and implementation of "Typing Tutor," a modular software system engineered from the ground up to simulate core operating system services. The project demonstrates a zero-dependency architecture by replacing standard C libraries (string.h, math.h, and the default malloc/free) with five custom-built foundational libraries. By integrating a low-level C backend with a reactive web-based frontend, the system achieves sub-millisecond precision in performance analytics, including Words Per Minute (WPM) and accuracy tracking. The research emphasizes custom memory allocation strategies, non-blocking I/O management, and session-based process state simulation.
\end{abstract}

\begin{IEEEkeywords}
Operating Systems, Custom Allocators, Systems Programming, Modular Architecture, Real-time I/O.
\end{IEEEkeywords}

\section{Introduction}
Modern software development often abstracts hardware complexities through high-level standard libraries. However, understanding the underlying system mechanics is critical for performance-critical applications. This project, "Typing Tutor," serves as a capstone exploration into low-level systems programming. By adhering to strict constraints—specifically the prohibition of standard library core logic—the system implements its own memory management, string manipulation, and I/O polling mechanisms. The final application provides an interactive terminal-style typing environment bridged to a modern web interface, offering a comprehensive demonstration of full-stack system integration.

\section{Functionality and System Modules}

\subsection{Memory Management}
The memory management module (\texttt{memory.c}) is the cornerstone of the system. It implements a custom first-fit allocator operating over a static 65KB virtual RAM region. Key features include:
\begin{itemize}
    \item \textbf{Ownership Tracking}: Blocks are tagged with \texttt{MEMORY\_OWNER\_SESSION} or \texttt{MEMORY\_OWNER\_PERSISTENT}, allowing the system to reset session-specific memory without affecting global datasets.
    \item \textbf{Block Metadata}: Each allocation is tracked via a descriptor table that manages offsets, sizes, and fragmentation status.
    \item \textbf{Zero Fragmentation Strategy}: A custom merging algorithm (\texttt{memory\_merge\_free\_blocks}) scans adjacent free blocks to maximize contiguous space.
\end{itemize}

\subsection{I/O Management}
Real-time interaction is handled by the \texttt{input\_engine.c} and \texttt{keyboard.c} modules. Unlike standard blocking input methods, the system utilizes:
\begin{itemize}
    \item \textbf{Non-blocking Polling}: The keyboard is configured via \texttt{fcntl} to operate in non-blocking mode, preventing the system from halting during user pauses.
    \item \textbf{Circular Buffer Queue}: Input is pushed into a custom circular buffer, allowing the backend to handle bursts of keystrokes with zero data loss.
\end{itemize}

\subsection{Process and Session Management}
While the application operates in a single address space, it simulates process-like state management through session control:
\begin{itemize}
    \item \textbf{Session Lifecycle}: Every typing test is treated as a new process context. \texttt{my\_reset()} ensures that the stack and heap-like structures are cleared before a new prompt is generated.
    \item \textbf{Halt and Error Handling}: The system includes a centralized error reporting mechanism in the \texttt{json\_loader}, preventing segmentation faults by validating dataset integrity at runtime.
\end{itemize}

\subsection{File System Interaction}
The system interacts with a simulated file-layer through the \texttt{json\_loader.c} module. It utilizes low-level system calls (\texttt{open}, \texttt{read}, \texttt{lseek}) to load a 650-word JSON dataset into persistent memory, simulating the behavior of a basic read-only file system driver.

\section{Methodology}
The development was split into two primary phases:
\subsubsection{Phase 1: Foundational Library Integration}
Construction of the five core libraries—\texttt{math.c}, \texttt{string.c}, \texttt{memory.c}, \texttt{screen.c}, and \texttt{keyboard.c}. These were tested for interoperability by creating a pipeline where keyboard input was parsed by strings and allocated via custom memory logic.
\subsubsection{Phase 2: System Integration and Frontend Bridging}
The C backend was integrated with a TCP socket layer (\texttt{server.c}) to communicate with a Vanilla JavaScript frontend. This phase involved implementing the dynamic sentence generation logic, which utilizes the custom math library to ensure fair word distribution and metrics calculation.

\section{Future Scope}
Future iterations of Typing Tutor aim to include:
\begin{itemize}
    \item \textbf{Multi-threaded Prompt Synthesis}: Using Pthreads to generate complex prompts in the background without affecting UI responsiveness.
    \item \textbf{Encrypted Leaderboards}: Implementing a custom XOR-based encryption module for secure local score storage.
    \item \textbf{TUI Interface}: A full terminal-based dashboard using raw ANSI escape codes for cursor positioning and color management.
\end{itemize}

\section{Conclusion}
The Typing Tutor project successfully demonstrates that high-level interactive applications can be built with complete control over system resources. By bypassing standard libraries, the project proves the viability of custom modular engines in educational and performance-critical contexts.

\begin{thebibliography}{00}
\bibitem{b1} A. Silberschatz, P. B. Galvin, and G. Gagne, \textit{Operating System Concepts}, 10th ed. Wiley, 2018.
\bibitem{b2} B. W. Kernighan and D. M. Ritchie, \textit{The C Programming Language}, 2nd ed. Prentice Hall, 1988.
\bibitem{b3} IEEE, "IEEE Standard for Information Technology—Portable Operating System Interface (POSIX)," IEEE Std 1003.1-2017.
\end{thebibliography}

\end{document}
